\documentclass[12pt,a4paper]{article}
\usepackage[utf8]{inputenc}
\usepackage[french]{babel}
\usepackage[T1]{fontenc}
\usepackage[margin=4cm]{geometry}
\usepackage{graphicx}
\usepackage{makeidx}
\usepackage{url}
\usepackage{caption}
\usepackage{subcaption}
\usepackage{enumerate}
\usepackage{enumitem}
\usepackage{ifthen}
\usepackage{amssymb}
\usepackage{dsfont}
\usepackage{multicol}
\usepackage[hidelinks]{hyperref}
\usepackage{amsmath}
\usepackage{multirow}
\usepackage{pgfplotstable}
\usepackage{booktabs}
\usepackage{csvsimple}
\usepackage[toc,page]{appendix} 

\newcommand{\figref}[1]{Figure~\ref{fig:#1}}
\newcommand{\tabref}[1]{Table~\ref{tab:#1}}
\newcommand{\secref}[1]{Section~\ref{sec:#1}}
\newcommand{\alref}[1]{Algorithm~\ref{alg:#1}}
\newcommand{\lstref}[1]{Listing~\ref{lst:#1}}
%\renewcommand{\appendixtocname}{Annexes} 
%\renewcommand{\appendixpagename}{Annexes} 


\title{\vspace{4cm}\textbf{L'impact du renommage sur les métriques de procédés}\vspace{3cm}}

\author{
Pierre Chanson\\\\
Mémoire de stage de Master2\\
Encadrants: Jean-Rémy Falleri et Matthieu Foucault\\\\
LaBRI, UMR 5800\\
F-33400, Talence, France\\\\
\texttt{Email: pierre.chanson@etu.u-bordeaux.fr,}\and
\texttt{\{falleri,mfoucault\}@labri.fr}\\
}


\begin{document}

\begin{figure}[t]
\center
\includegraphics[scale=0.25]{data/figures/UnivBordeaux.jpg}
\end{figure}
\maketitle
\tableofcontents
\newpage
\section{Introduction}

Process metrics such as Churn, which measures how much a Software artifact has been modified during its life cycle, or Number of developers, which measures how many developers did contribute to a software artifact, reveal the distribution of the workload of a Software project during its life cycle. They are defined to measure developer's activity with the main insight that developers habits have a high impact on the software quality~\cite{rahman_how_2013}. For instance, Bird et al. have shown that the Ownership metrics, which is a variant of the number of developers, has an impact on the number of post release failures in industrial software projects~\cite{bird_dont_2011}.

Measuring process metrics requires to consider a history of a Software project, contrary to classic metrics, such as LOC (Line of Code) or DIT (Depth of Inheritance Tree), that require just one snapshot of the project (usually the last one). Moreover, as software projects are nowadays managed by VCSs (Version Configuration Systems), the history must be composed of successive versions stored by VCS, and must reflect the real evolution of software artifacts. Further, a particular attention should be paid to artifact renaming that have been performed between versions as they may have a deep impact on the measure, and therefore may be an important threat to validity. For example, if we imagine a Software project where all the artifacts have been renamed in the middle of its life cycle, then needless to say that not considering them will return skewed measure for Churn or Number of developers metrics. 

In theory, artifact renaming does have an impact of process metrics. In practice, how many renaming occurs? When do they occur? And, what is their impact on process metrics? This paper aims to answer these questions. To that extent, it proposes a deep study of five real open source software projects with the objective to first propose descriptive statistics helping to understand the importance of artifact renaming in Software projects. The results we observed are highly important as renaming reach almost 50\% of the software artifacts for some projects. We then have computed two measures for process metrics, one that considers artifact renaming and one that does not, with the intent to evaluate the impact of renaming. Finally based on our observations, we propose some guidelines that aim at minimizing the threat of artifact renaming. For instance, we observed that there are less renaming in maintenance branches stored in VCS. As a consequence, when such a branch exists, it is better to compute process metrics on it.

This paper therefore provides the following contributions:
\begin{itemize}
\item A study of the artifact renaming phenomena made on five large open source Software projects. 
\item An impact analysis for process metrics. 
\item A set of guidelines aiming at reducing the threat of artifact renaming. 
\end{itemize}

This paper is structured as follows: Section~\ref{sec:rw} explains how process metrics are computed and presents how the artifact renaming has an impact on them. Section~\ref{sec:methodology} presents the detailed methodology of our study, including the construction of the dataset and the different analyses we realized. Section~\ref{sec:results} presents the main results of our study which shows that the artifact renaming phenomena does occur and does have an impact on process metrics. This section also presents our guideline to limit the threat of process metrics.
Finally, Section~\ref{sec:threats} of our study and Section~\ref{sec:conclusion} concludes this paper.






\section{Métriques de procédés}
\label{sec:metriques}

Nous décrivons dans cette partie comment les métriques de procédés sont calculés. Puis, nous expliquons comment les fichiers renommés peuvent avoir un impact sur ces métriques. Enfin, nous décrivons comment les VCS, actuellement, traitent le renommage.\\

\subsection{Calcul des métriques de procédés}

Les métriques de procédés (change metrics) permettent de calculer à quel point une entité de code source a été modifiée au cours d'une période donnée dans l'histoire d'un logiciel. Pour la prédiction de bugs, elles sont utilisés usuellement dans la dernière période, c'est à dire entre l'avant dernière et la dernière version du projet. L'objectif est de prédire les bugs qui apparaîtront lors de la prochaine version, en particulié si cette version est une \texttt{release}, c'est à dire une sortie de logiciel dite stable. Elles ne considèrent donc que les entités étant toujours présentes à la fin de la période et qui ont étés actives dans la période.\\

Un gestionnaire de versions (VCS) offre plusieurs moyens de calculer les métriques de procédés car il stocke les informations sur les entités modifiées à chaque nouvelle version, l'auteur de ces modifications, la date, etc. De plus, il permet la récupération du contenu de chaque entité et de l'ensemble d'un projet à une version donnée. Pour calculer ces métriques, il est donc possible d'analyser chaque entité modifiée lors d'une période puis de garder uniquement les entités toujours présentes à la dernière version de notre période.\\

Plus précisement, voici comment un VCS peut être utilisé pour calculer les trois métriques de procédés les plus utilisés: le nombre de développeurs (Number of Developers, NoD), le nombre de modifications  (Number of Changes, NoC) et le Code Churn (CC). 

\begin{enumerate}
\item  Nous récupérons d'abord la dernière version du projet pour obtenir les entités existantes à la fin de la période considérée. On note $A$ cet ensemble d'entités.
\item Toujours grâce aux commandes git, nous récupérons toutes les modifications effectuées durant la période. On note $C$ l'ensemble des modifications dans l'ordre chronologique.
\item  Troisièmement, nous parcourons cet ensemble de modifications en commençant par la plus ancienne ($c_0 \in C$) jusqu'à la plus récente ($c_n \in C$) dans le but de calculer les métriques de procédés pour chaque entité.
\end{enumerate}

 On note $\mu_{a}^{M}$ la valeur de la métrique $M$ pour l'entité $a$. Ensuite nous calculons les métriques comme suivant (on note $c_i$ la modification courante lors du parcours):

\begin{description}
	\item[NoD] (nombre de développeurs) Pour chaque entité $a$ pointé par $c_i$ qui appartient aussi à $A$ ($a \in A$), on ajoute à $\mu_{a}^{NoD}$ le nombre d'auteurs qui ont effectué les modifications $c_i$ et qui ont modifiés $a$ pour la première fois dans la période.
	\item[NoC] (nombre de modifications) Pour chaque entité $a$ pointé par $c_i$ qui appartient aussi à $A$ ($a \in A$), on ajoute $1$ à $\mu_{a}^{C}$ tels que $c_i$ indique qu'une nouvelle modification a été effectuée.
	\item[CC] (Code Churn) Pour chaque entité $a$ pointé par $c_i$ qui appartient aussi à $A$ ($a \in A$), on vérifie d'abord que la modification n'est pas une creation d'entité. Si c'est le cas celà veut dire que l'entité a été créée durant la période, donc on initialise son $\mu_{a}^{CC}$ à son nombre de lignes. Ensuite au prochain $c_j$ qui cible $a$ dans la période avec ($i < j$), on compare les deux versions et on ajoute à $\mu_{a}^{CC}$ le nombre de lignes ajoutées ou suprimées.
\end{description}




\subsection{Métriques de procédés et renommage}

Par ailleurs, il faut noter qu'un VCS identifie une entité par son chemin $+$ nom de fichier. On en déduit qu'un renommage du fichier ou d'un dossier, aura un impact sur le calcul des métriques. Pour expliquer cet impact, il est présenté un exemple d'historique d'un logiciel figure 2. Ce projet ne contient qu'une entité, Test.php, qui est renommé en Hello.php dans la dernière version. Dans cet exemple nous calculons le nombre de développeurs (NoD) entre la version 1 et 3.

Le NoD d'une entité de code source au cours d'une période de son histoire correspond au nombre de développeurs ayant été identifiés comme auteurs d'une modification sur l'entité pendant la période donnée.\\

\begin{figure}[t]
	\centering
	\includegraphics[width=0.8\linewidth,keepaspectratio]{data/figures/example.pdf}
	\caption{Exemple d'un historique de projet. Le projet est composé d'un seul fichier \texttt{Test.php} qui est renommé en \texttt{Hello.php} dans la dernière version.}
	\label{fig:example}
\end{figure}
En ne prenant pas en compte le renommage, la dernière version ne contient qu'une entité. C'est donc cette entité uniquement qui sera considérée. De plus, l'identité exacte de cette entité n'apparaît que lors de la version 3. Le calcul des métriques est donc trivial, NoD$ = 1$.

Par ailleurs, en prenant en compte le fait que ce fichier a été renommé, il y a trois versions à considérer en ce qui concerne l'entité. Le premier nom du fichier était Test.php. Ce fichier a un premier auteur lors de la version 1 puis un deuxième à la version 2. Le fichier est ensuite renommé en Hello.php par un troisième auteur. Le NoD est donc de 3.\\


Intéressons-nous aux outils disponibles pour la gestion de code source. Il existe un certain nombre de gestionnaires de versions tel que SVN, CVS, Mercurial ou Git qui pourraient être compatibles avec notre étude étant donné que nous avons simplement besoin de versions, c'est-à-dire un état du projet à un moment donné de son histoire, à comparer entre elles. Nous avons néanmoins étudié les VCS en détail et découvert que tous ne gèrent pas le renommage de fichiers de la même manière. La ~\tabref{vcs} résume notre étude. Alors que CVS ne gère pas du tout le renommage, SVN ou Mercural propose un mécanisme manuel de détection de renommage de fichiers. Git quant à lui propose un algorithme de détection de renommage automatique et optionnel. Pour les VCS qui utilisent une détection manuelle, cela implique que c'est aux développeurs d'utiliser les commandes appropriées. Cependant certaines études montrent que les développeurs n'utilisent pas ces commandes systématiquement. Le renommage peut être effectué jusqu'à $89\%$ du temps sans utiliser les commandes adaptées ~\cite{lavoie_inferring_2012,steidl_incremental_2014}. De plus l'étude de Kim et al montre que $51$\% des développeurs n'utilisent pas les commandes prévues par le VCS pour le refactoring (incluant le renommage). Ces trois études effectuées sur des projets open-source et industriels, montrent qu'il est dangereux de compter sur le fait que les développeurs utilisent les commandes adéquates pour le refactoring.\\


\subsection{Renommage et VCS}

\begin{table}[h]
\centering
\begin{tabular}{rccc}
\toprule
 & \multicolumn{3}{c}{Renaming handling}\\
\cmidrule{2-4}
& & \multicolumn{2}{c}{Automatic}\\
\cmidrule{3-4}
Tool & Manual & Standard & Optional\\
\midrule
CVS & & &\\
Subversion & $\times$ & &\\
Mercurial & $\times$ & & $\times$\\
Git & & & $\times$\\
\bottomrule
\end{tabular}
\caption{Traitement du renommage des principaux VCS.}
\label{tab:vcs}
\end{table}


\subsection{``Origin Analysis''}

Nous expliquons ici succinctement l'algorithme utilisé par Git pour la détection de renommage de fichiers. Celui-ci est connu sous le nom de ``Origin Analysis'' et est expliqué par Godfrey et al dans les articles ~\cite{tu_integrated_2002,godfrey_tracking_2002,godfrey_using_2005}.

Tout d'abord, il faut considérer deux versions successives d'un projet. Deux ensembles d'entités (fichiers, fonctions..) qui composent leurs versions respectives. Certaines entités ayant été modifiées, certaines supprimées et d'autres ajoutées.
La première analyse est une analyse de Bertillonage qui consiste à choisir un nombre de métriques, puis comparer les entités avec ces métriques. On compare alors les entités supprimées avec les entités ajoutées d'une version à l'autre. Grace à la distance Euclidienne calculée à partir des métriques combinés avec une comparaison des noms des entités, nous obtenons une liste des renommages potentiels.

Les analyses suivantes expliquées par Godfrey sont des améliorations de la première analyse, mais qui ne sont efficaces qu'à un niveau de granularité plus bas, c'est à dire une finesse d'analyse plus précise, en l'occurence au niveau des fonctions. Par exemple, l'analyse de dépendance qui traque les appels de fonctions, en comparant les fonctions appelantes et appelées. Ces analyses sont basées sur des seuils d'acceptabilité défini par l'utilisateur. Plus Godrey améliorera ces analyses, en prenant en compte par la suite les splits er merges de fonctions (algorithme inefficace au niveau des fichiers) plus l'utilisateur sera sollicité. (TODO: détailler ?)\\

\section{Méthodologie}
\label{sec:methodologie}

Nous présentons ici le déroulement de nos expérimentations qui consistent à étudier le phénomène du renommage et son impact sur le calcul des métriques de procédés. L’objectif étant d’évaluer si le renommage peut biaiser de manière significative les valeurs des métriques de procédés. Pour atteindre cet objectif, nous réalisons deux expérimentations successives.

Le but de la première expérience et de calculer la quantité de renommage durant les périodes de développement des logiciels. S’appuyant sur cette première expérience, notre deuxième expérience fournit une analyse de l’impact du renommage sur les métriques de procédés dans le pire des cas. 

Mais tout d'abord, nous avons besoin d'un ensemble de logiciels afin de former un coprus sur lequel appliquer nos expérimentations. 

\subsection{Corpus}

Nous avons sélectionné un ensemble de projets sur lesquels effectuer nos expérimentations. Des projets open-source, conséquents et connus de la communauté MSR. Ces projets qui sont présentés \tabref{projects} sont notamment régulièrement utilisés par l'équipe de Génie Logiciel au LaBRI. Ces $5$ projets forment un corpus comprenant différents langages de programmation ainsi qu'un nombre de lignes de code et un nombre de développeurs moyennement élevés à élevés par rapport à l'ensemble des projets open source utilisés couramment par la communauté. Les $5$ projets sont gérés sur Git afin de profiter de son mécanisme de traitement du renommage expliqué précédemment. \\

\begin{table*}[h]
\centering
\small
\begin{tabular}{rllp{1.7cm}c}
\toprule
Project & Main language & Size (LoC) & Number of developers & URL\\
\midrule
Jenkins & Java & 200851 & 454 & \url{github.com/jenkinsci/jenkins} \\
JQuery & JavaScript & 41656 & 223 & \url{github.com/jquery/jquery} \\
PHPUnit & PHP & 21799 & 152 & \url{github.com/sebastianbergmann/phpunit}\\
Pyramid & Python & 38726 & 205 & \url{github.com/Pylons/pyramid} \\
Rails & Ruby & 181002 & 2767 & \url{github.com/rails/rails}\\
\bottomrule
\end{tabular}
\caption{Notre corpus de projets.}
\label{tab:projects}
\end{table*}

Les projets de notre corpus suivent des phases distinctes durant leur cycle de vie. Habituellement une période de développement commence avant qu'une première sortie du logiciel, qu'on appelle \texttt{release}, soit accessible aux utilisateurs, puis cette release est maintenue pendant qu'une autre se prépare et ainsi de suite.

Nous avons ainsi deux phases, les phases de maintenance et les phases de développement. Nous divisons ces phases en périodes. Une période est délimité par deux releases et est composé d'un ensemble de versions successives. Chaque projet contient une période qui commence à la création du logiciel et qui se termine à la première release. On appelle cette période la période initiale. Les autres périodes peuvent être divisées en deux groupes, les périodes de releases majeures et les périodes de release mineures. On distingue les releases majeures des mineures par une augmentation significative du numéro de release. Par exemple $1.9-2.0$ pour Jquery, $3.7-4.0$ pour PHPUnit ou $0.13-1.0$ pour Rails. Usuellement, les périodes majeures contiennent un grand nombre de modifications en comparaison des périodes mineures.\\

Les conventions de nommage des releases sont spécifiques à chaque projet, par exemple:
\begin{itemize}
\item PHPunit: \texttt{3.5.0, 3.6.0} etc.
\item Pyramid: \texttt{1.0, 1.1} etc.
\item Jenkins: \texttt{jenkins-1\_400, jenkins-1\_410} etc.
\item Rails: \texttt{v2.0.0, v2.1.0} etc.
\end{itemize}
Nous avons analysé manuellement les dépôts des projets de notre corpus pour identifier leur releases mineures et majeures.\\

Le schéma \figref{model} permet de visualiser un dépôt Git ``type'', avec la branche master depuis laquelle, en remontant l'historique, on peut accéder au contenu de chaque release de la phase de développement et les branches de maintenance, dont la tête (la dernière version de la branche) n'est pas accessible depuis la branche master, qui composent la phase de maintenance.\\

\begin{figure}[h]
  \centering
  \includegraphics[scale=0.5]{data/figures/periods.pdf}
	\caption{Modèle d'architecture de dépot de code source}
	\label{fig:model}
\end{figure}

De plus, il faut noter que nous avons choisi d'exclure tous les fichiers qui ne sont pas du code source du corpus, étant donné que les métriques de procédés sont habituellement uniquement calculées sur ces fichiers. \\

\subsection{Première expérience}

 L'objectif de notre première expérience est de mieux comprendre le renommage d'entités. Nous souhaitons observer quand les renommages apparaissent et en quelle quantité. Nous avons donc analysé chaque période, comme décrites précédemment, sur chaque projet de notre corpus. 

Pour identifier les renommages, nous comptons sur le mécanisme de Git. Voici la procédure que nous avons suivie :
\begin{enumerate}
\item lister les fichiers existant à la fin de la période.
\item Pour chacun de ces fichiers, extraire sa séquence de modifications durant la période en activant la détection de renommage (commande \texttt{git log -M}).
\item Calculer à partir des informations recueillies le pourcentage de fichiers $\%F_{R}$ qui inclue au moins un renommage dans sa séquence de modification.
\end{enumerate}
\medskip

$\%F_{R}$, calculé sur une période représente en fait le pourcentage de \textbf{l'ensemble} des fichiers du projet qui ont été renommés dans cette période en particulier. $\%F_{R} = \frac{\#AF_{R}}{\#F}$ \\

Nous proposons en complément la même analyse que précédemment, mais en considérant uniquement les fichiers actifs dans cette même période. C'est à dire les fichiers qui ont été modifiés ou créés dans la période. On aura donc le pourcentage de fichiers \textbf{actifs} renommés dans la période : 
$\%AF_{R} = \frac{\#AF_{R}}{\#AF}$. \\
Etant donné que ${\#AF} <= {\#F}$, nous aurons toujours $\%F_{R} <= \%AF_{R}$. Cela nous permet de voir si les périodes qui contiennent le plus de renommage en interne sont aussi celles qui ont un impact sur l'ensemble du projet.\\

À notre connaissance, il n'existe pas d'évaluation empirique de l'algorithme utilisé par Git pour détecter les renommages. Néanmoins, nous procédons à une évaluation manuelle de son comportement dans la Section 4 et nous n'avons pas noté de faux positif sur 100 renommages aléatoires récupérés par notre analyse.\\

\subsection{Deuxième expérience}

L'objectif de la deuxième expérience est de voir si le renommage peut biaiser significativement les valeurs des métriques de procédés décrites ci-dessus. Pour ça, nous effectuons une analyse dans le pire des cas. Nous sélectionnons une période par projet, celle qui a la plus grande valeur de fichiers renommés en excluant la période initiale qui n'est généralement pas observée dans les études. Nous calculons ensuite les trois métriques avec et sans le renommage de fichiers pris en compte. Puis nous calculons la corrélation de coefficient de Spearman entre les métriques avec et sans la détection de renommage. Un coefficient élevé, proche de $1$, indiquera que les métriques avec et sans détection de renommage sont très similaires alors qu'un coefficient plus petit, 0.5 et moins, indiquera que les métriques avec et sans détection de renommage sont très différentes.\\  

\section{Résultats}
\label{sec:resultats}

\subsection{Première expérience}

Les résultats de la première expérience sont montrés dans la \figref{renaming}. Tout d'abord, le nombre de renommage varie beaucoup entre les projets. Par exemple, Jenkins a au plus $10\%$ de ses fichiers renommés dans la pire période alors que PHPUnit a deux périodes à plus de $50\%$. Le nombre de renommages varie aussi en fonction des périodes, par exemple dans PHPUnit la période $3.6 - 3.7$ a moins de $5\%$ de fichiers renommés alors que la période $3.7 - 4.0$ a presque $99\%$. En général, il y a beaucoup de périodes avec $0\%$ de fichiers renommés.\\

\begin{figure*}[h]
	\centering
	\includegraphics[width=0.85\linewidth,keepaspectratio]{data/figures/renaming.pdf}
	\caption{Pourcentage de fichiers renommés ($\%F_R$) dans chaque période de chaque projet de notre corpus. La période initiale est en gris foncé, les périodes majeures en gris et les périodes mineures en gris clair.}
	\label{fig:renaming}
\end{figure*}

Par rapport à la localisation de ces renommages, la période initiale semble la plus prolifique au renommage. En général, elle contient le plus grand nombre de fichiers renommés (sauf pour PHPUnit). Les périodes de développement sont plus susceptibles d'avoir des renommages que les périodes de maintenance. Ainsi, les $5$ projets sont quasiment à $0\%$ de fichiers renommés dans les périodes de maintenance. Finalement, certaines périodes de développement peuvent contenir beaucoup de renommages. Les résultats montrent que les releases majeures sont souvent les pires périodes de développement en nombre de fichiers renommés: C'est le cas pour PHPUnit et Rails alors que Jenkins et Pyramid ne contiennent pas de releases majeures.\\

%Le détail des résultats obtenus par notre outil de détection de renommage sont montrés dans les tables en annexes \tabref{jenkins}, \tabref{jquery}, \tabref{phpunit}, \tabref{pyramid} et %\tabref{rails}. Ces tables mettent en avant les valeurs suivantes:\\

%\begin{itemize}
%\item \emph{Nombre de fichiers ($\#F$):} nombre de fichiers dans le projet à la dernière version de la période..
%\item \emph{Nombre de fichiers actifs ($\#AF$):} nombre de fichiers créés, suprimés, copiés ou renommés durant la période et présents à la dernière version de la période.
%\item \emph{Pourcentage de fichiers actifs ($\%AF$):} $\%AF = \frac{\#AF}{\#F}$.
%\item \emph{Nombre de fichiers actifs renommés ($\#AF_{r}$):} nombre de fichiers actifs qui ont étés renommés.
%\item \emph{Pourcentage de fichiers renommés ($\%F_{R}$):} $\%F_{R} = \frac{\#AF_{R}}{\#F}$.
%\item \emph{Pourcentage de fichiers actifs renommés ($\%AF_{R}$):} $\%AF_{R} = \frac{\#AF_{R}}{\#AF}$.
%\end{itemize}

En ce qui concerne le pourcentage de fichiers \textbf{actif} renommés ($\%AF_R$) dans les périodes de nos projets, la \figref{renaming2} présente nos résultats. Nous pouvons remarquer que le taux de fichiers renommés ne change pas ou augmente légèrement pour Jenkins par exemple. Etant donné que $\%AF_R >= \%F_R$ comme nous l'avons expliqué plus tôt, on en déduit qu'une grande partie des périodes, en tout cas les périodes qui ont un impact sur l'ensemble du projet, accèdent à la totalité ou en grande partie (Jenkins) des fichiers du projet durant la période. \\

\begin{figure*}[h]
	\centering
	\includegraphics[width=0.85\linewidth,keepaspectratio]{data/figures/renaming2.pdf}
	\caption{Pourcentage de fichiers actifs renommés ($\%AF_R$) dans chaque période de chaque projet de notre corpus. La période initiale est en gris foncé, les périodes majeures en gris et les périodes mineures en gris clair.}
	\label{fig:renaming2}
\end{figure*}

\subsection{Deuxième expérience}
Les résultats de notre deuxième expérience sont montrés dans la \tabref{spearman}. Ils montrent que la corrélation de coefficients de Spearman entre les métriques de procédés avec et sans détection de renommages dépendent beaucoup de la période et de la métrique choisie. Les métriques de procédés ne sont pas affectées par le renommage dans les projets Jenkins, Rails et Pyramid. Ainsi, le coefficient de corrélation est proche de $1$ dans tous les cas. D'un autre côté, pour PHPUnit et JQuery les métriques peuvent être sévèrement impactées par le renommage. Pour JQuery, la métrique Code Churn n'est pas affectée par le renommage, mais NoD et NoC sont quant à eux significativement impactés. Pour PHPUnit, toutes les métriques sont affectées par le renommage. Sur ces deux derniers projets, la métrique la plus sensible aux renommages de fichiers est le nombre de développeurs (NoD). Sur ces deux derniers projets, la métrique la plus sensible aux renommages de fichiers est le nombre de développeurs (NoD).

Finalement, on peut noter que seules les périodes ayant eu un grand pourcentage de fichiers renommés ($\%F_R$) ont été impactées. On peut aussi noter que des métriques biaisées par le renommage ont été calculées dans des releases majeures et mineures.\\

Nous avons étudié manuellement les deux périodes qui ont affecté les valeurs des métriques de procédés (JQuery 1.1 - 1.2 et PHPUnit 3.7 - 4.0). Dans ces deux périodes, la structure générale du projet a été modifié. Notamment, des renommages de dossiers à la racine ont étés effectués. Par conséquent, un grand nombre de fichiers a été renommé de manière transitive. C'est une pratique courante dans le développement logiciel, donc le phénomène pourrait apparaitre dans n'importe quelle période ou projet. Il est intéressant de noter que dans ces deux périodes, les changements de structure ont été effectués en grande partie dans un seul commit proche de la fin de la période.\\


\begin{table*}[h]
\centering
\csvreader[tabular=rcccc, table head=\toprule & & \multicolumn{3}{c}{Change metrics}\\\cmidrule{3-5} Period & $\%F_R$ & CC & NoD & NoC\\\midrule, late after line=\\, late after last line=\\\bottomrule]{data/tables/correlations.csv}%
{1=\period,2=\fr,3=\churnall,4=\devall,5=\modificationsall}%
{\period & \fr & \churnall & \devall & \modificationsall}
\caption{La corrélation de coefficients de Spearman entre les valeurs des métriques de procédés avec et sans détection de renommage. Les codes de signification sont: *** $\leq 0.01$, ** $\leq 0.05$, * $\leq 0.1$ et ! $> 0.1$. Les coefficients moyen et faible sont affichés en gras.(TODO : gérer les tableaux)}
\label{tab:spearman}
\end{table*}


\subsection{Validations et limitations}

Notre étude fait l'hypothèse que les renommages détectés par Git sont corrects. Néanmoins, nous n'avons pas rencontré d'évaluation empirique de l'algorithme de détection de renommage de Git et par conséquent nous n'avons pas confiance dans ses résultats. Afin d'atténuer cette menace, nous avons tiré au hasard $100$ renommages détectés par Git lors de notre expérimentation. Nous avons évalué manuellement chaque renommage de fichier pour vérifier si la détection était correcte. Vérifier que la détection soit correcte consiste à s'assurer que le contenu du fichier avant et après renommage soit très similaire, et qu'aucun autre fichier ajouté dans la même version n'ait un contenu similaire. Cette analyse manuelle nous a révélé que  $100\%$ des renommages étaient correctement détectés. Même si nous sommes conscients que l'algorithme Git peut donner des faux positifs, cette expérience montre que l'utilisation de Git dans la détection de renommage reste raisonnable. Nous n'avons pas analysé les faux négatifs, c'est à dire les vrais renommages de fichiers non détectés par Git, car ça ne risquerait pas de fausser nos conclusions. Cela ne pourrait que diminuer les coefficients de corrélation si nous l'avons sous-estimé.\\  

En raison d'un très petit nombre de fichiers, les valeurs des coéfficient de corrélation pour JQuery sont très faible pour NoD et NoC. Ces coefficiant nous montrent que cette période est fortement affectée par le renommage. Cependant, le fait que les coefficiants de corrélation soient trop bas nous amène à une interprétation qui ne remet pas en cause nos conclusions.\\

Nous avons uniquement évalué l'effet du renommage sur trois métriques de procédés. L'effet du renommage pourrait être différent (mieux ou pire) pour les autres métriques de procédés. Notre solution pour limiter ce risque était d'analyser les métriques de procédés les plus utilisées comme il est montré dans ~~\cite{radjenovic_software_2013}. Les autres métriques de procédés sont souvent basées sur ces trois métriques comme la métrique \emph{code ownership}~\cite{bird_dont_2011} ou \emph{module activity focus}~\cite{posnett_dual_2013}.\\

Pour la métrique NoD, nous n'avons pas appliqué un algorithme de merge d'identités ~\cite{goeminne_comparison_2013}. Il pourrait en résulter des valeurs incorrectes. Cependant, ce phénomène est susceptible d'arriver pour les deux calculs de métrique avec et sans le renommage, donc le risque qu'il invalide nos conclusions est faible.\\

Concernant notre conclusion sur la quantité de renommage, le corpus utilisé ne garantit pas qu'elle puisse être généralisée. En effet, nous avons uniquement utilisé des projets open-source, alors que les projets industriels sont connus pour être sensiblement différents. En ce qui concerne la validité sur les projets open-source, notre corpus est trop petit pour généraliser cette conclusion.\\

\section{Analyse des études antérieures et recommandations}
\label{sec:analyse}

Dans cette partie, nous procédons à l’analyse d’études antérieures sur la prédiction de bugs qui ont utilisé les métriques de procédés pour la prédiction de bugs. Nous évaluons si les valeurs des métriques de procédés pourraient être biaisées en regardant la façon dont ils ont recueilli leurs données.

Enfin, nous donnons quelques lignes de conduite à suivre pour aider les chercheurs et développeurs à éviter l’impact que le renommage d’entités pourrait avoir sur les métriques de procédés.  

\subsection{Analyse des études antérieures}

Premièrement, comme nous l'avons montré dans la \secref{resultats}, il est important de remarquer que les périodes contenant un taux de fichiers renommés élevés sont rares. Ainsi, la majeure partie des études antérieures ne devraient pas être affectées par ce phénomène. 

De plus, même dans le cas d'analyses sur des périodes contenant un taux élevé de fichiers renommés, les résultats de ces analyses pourraient également être améliorés, car les métriques de procédés auraient probablement été sous-estimées. Toutefois, plusieurs études antérieures peuvent être affectées par le renommage, comme nous le signalerons dans la suite de cette section. Quantifier de tels effets sur les études antérieures est hors de portée de notre sujet, mais nous fournissons tout de même certaines lignes de conduite à respecter pour les études futures.\\

Dans notre examen des études antérieures, nous n'analysons que $26$ des articles référencés dans ~\cite{radjenovic_software_2013}. Les articles qui utilisent les trois métriques de procédés, CC, NoD ou NoC. Cependant, certaines des autres études référencées dans cet article utilisent d'autres métriques de procédés et pourraient donc aussi être affectées par le renommage.\\

$15$ de ces études analysent des projets industriels, ~\cite{arisholm_systematic_2010,graves_predicting_2000,khoshgoftaar_using_2000,layman_iterative_2008,munson_code_1998,nagappan_use_2005,nagappan_influence_2008,nagappan_using_2007,nagappan_using_2006,nagappan_change_2010,nikora_building_2006,ostrand_programmer-based_2010,weyuker_too_2008,weyuker_using_2007,yuan_application_2000}. Aucune de ces études ne parle de renommage. Mais le manque d'informations récoltées sur les VCS utilisés et sur le projet en lui-même, ne nous permet pas de savoir si le renommage pouvait avoir un impact sur ces projets. Néanmoins, l'article de Kim et al ~\cite{kim_field_2012} explique que les développeurs dans son étude effectuent des opérations de refactoring, dont du renommage, sans utiliser les outils du VCS appropriés. Ainsi, ces études pourraient être impactées par le renommage en fonction des outils utilisés et des habitudes de développement.\\

$11$ études analysent des logiciels open-source \cite{dambros_relationship_2009,bacchelli_are_2010,caglayan_merits_2009,dambros_evaluating_2012,dambros_evaluating_2012,dambros_extensive_2010,illes-seifert_exploring_2010,li_finding_2005,matsumoto_analysis_2010,moser_analysis_2008,moser_comparative_2008,schroter_if_2006}. Les VCS utilisés dans ces études sont CVS ou Subversion. CVS ne gère pas le renommage et Subversion uniquement de manière manuelle ce qui est dangereux comme expliqué dans l'article ~\cite{lavoie_inferring_2012,steidl_incremental_2014}. Seulement deux de ces études ~\cite{moser_analysis_2008,moser_comparative_2008} parlent de renommage, dans leur set de données ou dans les ''Threats to validity''. Pour réduire le risque d'erreur dans leurs expérimentations, ces deux études ont supprimé systématiquement tous les fichiers ajoutés ou supprimés durant les périodes analysées. C'est un bon moyen d'éviter de calculer des métriques de procédés biaisés, mais cela implique aussi de supprimer inutilement, du jeu de données, un nombre significatif de fichiers.\\

\subsection{Recommandations}
\label{sec:guidelines}

Les résultats de nos deux expérimentations nous permettent de déduire de simples lignes de conduite pour calculer les métriques de procédés. Voici donc nos recommandations: 
\begin{itemize}
\item Eviter de calculer ces métriques durant les périodes initiales. En effet, ces périodes contiennent habituellement une quantité de renommage significative. Comme nous l'avons vu, les deux périodes majeures et mineures peuvent contenir un beaucoup de renommage, bien que les releases majeures semblent plus sujettes au renommage. 
\item Utiliser systématiquement un algorithme de détection de renommage, afin d'éviter d'analyser les mauvaises périodes. Git propose un algorithme dédié qui semble avoir une bonne précision mais un rappel inconnu. Par conséquent, l'utilisation de projets gérés avec Git parait la méthode la plus simple pour diminuer les risques du renommage. Des algorithmes de détection plus avancés sont décrits dans la littérature ~\cite{antoniol_automatic_2004,lavoie_inferring_2012,steidl_incremental_2014}. Ils ont été validés par des études empiriques donc ils pourraient réaliser une analyse meilleure que celle de Git. 
\item Pour les métriques de procédés calculées à des niveaux de granularité plus fins que celui des fichiers, utiliser les algorithmes d``Origin Analysis'' tels que ~\cite{wu_aura:_2010}. Ces algorithmes sont efficaces au niveau de granularité des fonctions.
\item Le renommage d'entités peut être un risque important, penser à indiquer systématiquement comment il a été traité dans les études futures.\\       
\end{itemize}

\section{Conclusion}
\label{sec:conclusion}

Dans cet article, nous avons évalué l'impact du renommage d'entités sur les valeurs des métriques de procédés logicielles. Nous avons effectué une étude empirique sur cinq projets open-source connus et matures. Nous avons observé que les périodes initiales des projets sont plus enclines à contenir du renommage que les autres périodes. Plus important, nous avons constaté que d'autres périodes peuvent contenir une quantité importante de renommage, en particulier celles correspondantes à la mise au point de releases majeures. Enfin, nous avons observé que le renommage pouvait biaiser considérablement les valeurs des métriques de procédés. Par conséquent, les chercheurs et développeurs devraient être prudent lors du calcul des métriques de procédés. Nous recommandons d'éviter le calcul des métriques de procédés lors des périodes initiales. Pour finir, nous recommandons fortement d'utiliser un algorithme de détection de renommage lors du calcul des métriques de procédés sur d'autres périodes tels qu'il pourrait en biaiser fortement le résultat.\\

A l'avenir, nous prévoyons d'évaluer la précision des algorithmes existants de détection de renommage. Nous prévoyons également d'évaluer l'impact de la fusion de code (code merging) sur les métriques de procédés.

\newpage
\bibliographystyle{plain}
\bibliography{renaming}
%\newpage
%\begin{appendices}
%\label{sec:tables}

\begin{table}[h]
\centering
\scriptsize
\csvreader[tabular=rcccccc, table head=\toprule & & \multicolumn{5}{c}{Renaming metrics}\\\cmidrule{3-7} Period & Kind & $\#F$ & $\#AF$ & $\%AF$ & $\%F_r$ & $\%AF_r$\\\midrule, late after line=\\, late after last line=\\\bottomrule, after table=]{data/tables/jenkins.csv}%
{1=\period,2=\kind,3=\nf,4=\naf,5=\paf,6=\pfr,7=\pafr}%
{\period & \kind & \nf & \naf & \paf & \pfr & \pafr}
\caption{Amount and location of renaming in Jenkins}
\label{tab:jenkins}
\end{table}

\begin{table}[h]
\centering
\scriptsize
\csvreader[tabular=rcccccc, table head=\toprule & & \multicolumn{5}{c}{Renaming metrics}\\\cmidrule{3-7} Period & Kind & $\#F$ & $\#AF$ & $\%AF$ & $\%F_r$ & $\%AF_r$\\\midrule, late after line=\\, late after last line=\\\bottomrule]{data/tables/jquery.csv}%
{1=\period,2=\kind,3=\nf,4=\naf,5=\paf,6=\pfr,7=\pafr}%
{\period & \kind & \nf & \naf & \paf & \pfr & \pafr}
\caption{Amount and location of renaming in JQuery}
\label{tab:jquery}
\end{table}

\begin{table}[h]
\centering
\scriptsize
\csvreader[tabular=rcccccc, table head=\toprule & & \multicolumn{5}{c}{Renaming metrics}\\\cmidrule{3-7} Period & Kind & $\#F$ & $\#AF$ & $\%AF$ & $\%F_r$ & $\%AF_r$\\\midrule, late after line=\\, late after last line=\\\bottomrule]{data/tables/phpunit.csv}%
{1=\period,2=\kind,3=\nf,4=\naf,5=\paf,6=\pfr,7=\pafr}%
{\period & \kind & \nf & \naf & \paf & \pfr & \pafr}
\caption{Amount and location of renaming in PHPUnit}
\label{tab:phpunit}
\end{table}

\begin{table}[h]
\centering
\scriptsize
\csvreader[tabular=rcccccc, table head=\toprule & & \multicolumn{5}{c}{Renaming metrics}\\\cmidrule{3-7} Period & Kind & $\#F$ & $\#AF$ & $\%AF$ & $\%F_r$ & $\%AF_r$\\\midrule, late after line=\\, late after last line=\\\bottomrule]{data/tables/pyramid.csv}%
{1=\period,2=\kind,3=\nf,4=\naf,5=\paf,6=\pfr,7=\pafr}%
{\period & \kind & \nf & \naf & \paf & \pfr & \pafr}
\caption{Amount and location of renaming in Pyramid}
\label{tab:pyramid}
\end{table}

\begin{table}[h]
\centering
\scriptsize
\csvreader[tabular=rcccccc, table head=\toprule & & \multicolumn{5}{c}{Renaming metrics}\\\cmidrule{3-7} Period & Kind & $\#F$ & $\#AF$ & $\%AF$ & $\%F_r$ & $\%AF_r$\\\midrule, late after line=\\, late after last line=\\\bottomrule]{data/tables/rails.csv}%
{1=\period,2=\kind,3=\nf,4=\naf,5=\paf,6=\pfr,7=\pafr}%
{\period & \kind & \nf & \naf & \paf & \pfr & \pafr}
\caption{Amount and location of renaming in Rails}
\label{tab:rails}
\end{table}

%\end{appendices} 

\end{document}
